\documentclass[letterpaper,11pt]{article}
\usepackage[T1]{fontenc}
\usepackage[utf8]{inputenc}
\usepackage{lmodern}
\usepackage[french]{babel}
\usepackage[margin=1in,headheight=28pt]{geometry}
\renewcommand{\familydefault}{\sfdefault}
\usepackage{amsmath,amssymb,amsthm,bm,enumitem}
\usepackage[dvipsnames]{xcolor}
\usepackage[small,compact]{titlesec}
\usepackage{fancyhdr,fancyvrb}
\usepackage[colorlinks=true,allcolors=blue]{hyperref}
\pagestyle{fancy}
\fancyhead[L]{STT-2100 \\ Régression}
\fancyhead[R]{A.-A. Vallée \\ Université Laval}
\fancyfoot[C]{\thepage}
\setlength{\parindent}{0pt}
\setlength{\parskip}{5pt}
\fvset{fontsize=\small,frame=single,framesep=2mm}
\newcommand{\esp}{\mathbb{E}}
\newcommand{\var}{\operatorname{Var}}
\newcommand{\bo}[1]{\bm{#1}}
\newcommand{\N}{\mathcal{N}}
% Définir \avecsolutions avant l'input pour une version corrigée.
\newcommand{\prof}[1]{\ifdefined\avecsolutions\par{\color{blue}\textbf{Solution :} #1}\par\fi}
\begin{document}
\begin{center}{\Large Observations aberrantes et influentes avec R}\end{center}
Cette fiche accompagne les notes du chapitre 10. Les noms des objets peuvent être adaptés à vos données. Les exemples utilisent des données incluses dans R ou des données simulées; aucun fichier externe n’est requis.

\par\medskip\noindent\begin{minipage}{\linewidth}
\subsection*{Calculer les diagnostics}
Un grand levier concerne le profil explicatif; un grand résidu concerne l'écart de réponse au modèle; la distance de Cook mesure l'effet sur l'ajustement. Les seuils ci-dessous signalent des observations à examiner.
\begin{Verbatim}
donnees <- mtcars
mod <- lm(mpg ~ wt + hp, data = donnees)
n <- nobs(mod)
pprime <- length(coef(mod))
diagnostic <- data.frame(
  voiture = rownames(model.frame(mod)),
  levier = hatvalues(mod),
  standardise = rstandard(mod),
  studentise = rstudent(mod),
  cook = cooks.distance(mod),
  dffits = dffits(mod),
  covratio = covratio(mod))
diagnostic$signal <- with(diagnostic,
  levier > 2 * pprime / n | abs(studentise) > 2 |
  cook > 4 / (n - pprime))
diagnostic[order(diagnostic$cook, decreasing = TRUE), ]
\end{Verbatim}

\end{minipage}\par
\par\medskip\noindent\begin{minipage}{\linewidth}
\subsection*{Changements des coefficients}
dfbeta() donne les différences non standardisées; dfbetas() donne les différences standardisées. Ces quantités ne sont pas interchangeables.
\begin{Verbatim}
dfbeta(mod)
dfbetas(mod)
influence.measures(mod)
par(mfrow = c(1, 2))
plot(mod, which = 4)
plot(mod, which = 5)
par(mfrow = c(1, 1))
\end{Verbatim}

\end{minipage}\par
\par\medskip\noindent\begin{minipage}{\linewidth}
\subsection*{Analyse de sensibilité}
Le retrait ci-dessous sert à mesurer l'effet de l'observation ayant la plus grande distance de Cook. Vérifiez la donnée et le contexte avant toute décision de retrait définitif.
\begin{Verbatim}
i <- which.max(cooks.distance(mod))
rownames(donnees)[i]
mod_sans <- lm(mpg ~ wt + hp, data = donnees[-i, ])
cbind(avec = coef(mod), sans = coef(mod_sans))
profil <- data.frame(wt = 3, hp = 110)
rbind(avec = predict(mod, newdata = profil),
      sans = predict(mod_sans, newdata = profil))
\end{Verbatim}

\end{minipage}\par
\par\medskip\noindent\begin{minipage}{\linewidth}
\subsection*{Dépistage simultané des valeurs aberrantes}
La correction de Bonferroni contrôle le seuil familial sous les postulats du modèle. Le résidu studentisé externe utilise une variance estimée après retrait, avec n-pprime-1 degrés de liberté.
\begin{Verbatim}
alpha <- 0.05
seuil <- qt(1 - alpha / (2 * n), df = n - pprime - 1)
diagnostic[abs(diagnostic$studentise) > seuil, ]
\end{Verbatim}

\end{minipage}\par
\par\medskip\noindent\begin{minipage}{\linewidth}
\subsection*{Aide}
Dans R, utilisez \texttt{?lm} et l’aide des fonctions présentées. Voir aussi \href{https://stat.ethz.ch/R-manual/R-devel/library/stats/html/influence.measures.html}{la documentation des diagnostics d’influence}.
\end{minipage}\par
\end{document}
